\documentclass[11pt]{article}
\usepackage[utf8]{inputenc}
\usepackage[margin=1in]{geometry}
\usepackage{booktabs}
\usepackage{graphicx}
\usepackage{hyperref}
\usepackage{listings}
\usepackage{xcolor}
\usepackage{enumitem}

\definecolor{codebg}{RGB}{245,245,245}
\definecolor{codeframe}{RGB}{200,200,200}
\lstset{
  backgroundcolor=\color{codebg},
  frame=single,
  rulecolor=\color{codeframe},
  basicstyle=\ttfamily\small,
  breaklines=true,
  columns=fullflexible,
  language=Python,
  keywordstyle=\color{blue},
  commentstyle=\color{gray},
  stringstyle=\color{red!70!black},
}

\title{\textbf{Bonus: Ray Serve Integration} \\ \large An Additional Serving Framework for the GRU Ranker}
\author{Serving Team Member}
\date{April 2026}

\begin{document}
\maketitle

\section{Overview}

For the extra credit, I integrated \textbf{Ray Serve} as an additional serving framework beyond the lab-required FastAPI baseline and Triton Inference Server. Ray Serve is a scalable model serving library built on Ray that provides built-in replica management and application-level request batching --- capabilities that neither vanilla FastAPI nor Triton offer in the same way for CPU-bound, small-model workloads.

\medskip
\noindent\textbf{Artifacts submitted:}
\begin{itemize}[nosep]
  \item \texttt{ray\_serve\_api.py} --- Ray Serve deployment with batching
  \item \texttt{Dockerfile.ray} --- Container definition for Chameleon deployment
\end{itemize}

\section{Why Ray Serve Improves the Serving Design}

The lab frameworks each have limitations for our use case (a lightweight 314K-parameter GRU Ranker on CPU):

\begin{itemize}
  \item \textbf{FastAPI baseline} (lab default): Single-process, single-model serving. Scaling requires external tools (e.g., Gunicorn workers, Kubernetes HPA). No built-in request batching.
  \item \textbf{Triton Inference Server} (lab framework): Designed for GPU-heavy, large-model workloads. For our small CPU model, Triton's double-hop architecture (client $\to$ HTTP $\to$ Triton $\to$ inference $\to$ HTTP $\to$ client) introduces significant serialization overhead, resulting in \textbf{380\,ms p50 latency} --- the worst of all options tested.
\end{itemize}

\noindent\textbf{Ray Serve addresses both limitations:}
\begin{enumerate}
  \item \textbf{Built-in horizontal scaling:} The \texttt{@serve.deployment(num\_replicas=2)} decorator automatically manages multiple model replicas, distributing load without external orchestration.
  \item \textbf{Application-level batching:} The \texttt{@serve.batch(max\_batch\_size=8)} decorator transparently collects incoming requests and batches them before model inference, amortizing per-request overhead.
  \item \textbf{In-process execution:} Unlike Triton's proxy architecture, Ray Serve runs the model within the same process group, avoiding HTTP serialization penalties.
\end{enumerate}

\section{Implementation Details}

The core deployment is defined in \texttt{ray\_serve\_api.py}:

\begin{lstlisting}
@serve.deployment(
    num_replicas=2,
    ray_actor_options={"num_cpus": 1},
)
@serve.ingress(app)
class GRURankerDeployment:
    def __init__(self):
        self.model = load_or_create_model(ckpt, self.device)
        self.ctx = ServingContext(...)

    @serve.batch(max_batch_size=8, batch_wait_timeout_s=0.05)
    async def score_batch(self, batch_inputs):
        # Stack inputs across batch dimension
        # Run batched inference through GRU model
        # Return per-request score arrays

    @app.post("/predict")
    async def predict(self, req: RecommendRequest):
        # Preprocess: embedding lookup, candidate generation
        # Call self.score_batch(inp)  -- auto-batched
        # Return top-5 recommendations
\end{lstlisting}

\noindent Key design decisions:
\begin{itemize}[nosep]
  \item \textbf{2 replicas} with 1 CPU each, matching our 4-vCPU Chameleon instance (leaving headroom for the OS and preprocessing).
  \item \textbf{Batch timeout of 50\,ms}: Balances latency (requests are not delayed too long waiting for a full batch) with throughput (enough time to accumulate requests under load).
  \item \textbf{Async predict endpoint}: Enables non-blocking I/O while the batch collects, improving concurrency handling.
\end{itemize}

\section{Benchmark Results}

All experiments ran on Chameleon \texttt{compute\_haswell} (Intel Haswell, 4\,GB RAM), inside Docker containers, using \texttt{benchmark\_serving.py} with 200 requests per run.

\begin{table}[h]
\centering
\caption{Ray Serve vs.\ Lab Frameworks --- Performance Comparison}
\label{tab:rayserve_comparison}
\resizebox{\textwidth}{!}{%
\begin{tabular}{llrrrl}
\toprule
\textbf{Option} & \textbf{Framework} & \textbf{p50 (ms)} & \textbf{p95 (ms)} & \textbf{Throughput (req/s)} & \textbf{Notes} \\
\midrule
baseline\_pytorch (c=10) & FastAPI + PyTorch & 70.76 & 92.00 & 137.86 & Lab baseline \\
onnx\_int8 (c=10) & FastAPI + ONNX RT & 70.10 & 88.68 & 139.84 & Best model-level opt. \\
triton\_batched (c=10) & Triton Inf. Server & 380.33 & 520.75 & 25.88 & Lab system-level opt. \\
\midrule
\textbf{ray\_serve (c=10)} & \textbf{Ray Serve} & \textbf{133.82} & \textbf{192.95} & \textbf{71.07} & \textbf{Bonus framework} \\
\textbf{ray\_serve (c=50)} & \textbf{Ray Serve} & \textbf{632.78} & \textbf{773.08} & \textbf{77.78} & \textbf{Stress test (5$\times$ load)} \\
\bottomrule
\end{tabular}%
}
\end{table}

\section{Key Finding: Graceful Degradation Under Load}

The most significant result is Ray Serve's behavior under a \textbf{5$\times$ concurrency increase} (10 $\to$ 50 concurrent users):

\begin{itemize}
  \item \textbf{Throughput held steady:} 71.07 $\to$ 77.78 req/s (actually \emph{increased} slightly due to better batch utilization).
  \item \textbf{Zero errors:} 0\% error rate at both concurrency levels.
  \item \textbf{Latency scaled proportionally:} p50 increased from 134\,ms to 633\,ms, which is expected --- with 5$\times$ more concurrent users, each request waits longer in the queue, but no requests are dropped.
\end{itemize}

\noindent This demonstrates that Ray Serve's replica-based architecture with batching provides \textbf{effective horizontal scaling} that neither vanilla FastAPI (which would require external load balancing) nor Triton (which showed poor performance for this small model) can match out-of-the-box.

\medskip
\noindent\textbf{Comparison to Triton:} Ray Serve achieved \textbf{2.7$\times$ higher throughput} than Triton (71 vs.\ 26 req/s) at the same concurrency, because it avoids the HTTP serialization double-hop that dominates Triton's latency for small, CPU-bound models.

\section{Conclusion}

Ray Serve is a strong fit for serving lightweight CPU models like our GRU Ranker. Its built-in replica management and \texttt{@serve.batch} decorator provide scaling capabilities that go beyond the lab frameworks, with minimal code overhead. For production deployment, Ray Serve would pair well with the ONNX INT8 optimized model to combine the best model-level and system-level optimizations.

\end{document}
